\chapter{Adatbevitel: typed SQL, CRUD, formok és PRG}

\noindent\textbf{Tanulási út: Gyors út: alap.} Adj bounded formokat és CRUD-ot az alkalmazáshoz az authorization és transaction határok gyengítése nélkül.\par\medskip

\rustbridgeblockhu{A typed értékek, structok, \code{Option}/\code{Result}, összehasonlítások és normál vezérlés maradnak a business rule-ok nyelve.}{A route form deklaráció, validációs boundok, transaction capability, authorization proof és typed redirect definiálja a védett mutation-határt.}{Ne írd át a business validációt framework-mágiává. A domain logika maradjon Rust-szerű, Velran-elemet ott használj, ahol adat trust- vagy effect-határon lép át.}


\invariantblock{Biztonsági invariáns: input korlátozása effect előtt}{A request adatát kezeld nem megbízhatóként mindaddig, amíg a várt típusra nem parse-oltad, szükség szerint nem normalizáltad, és explicit boundokkal nem validáltad. Ez adatbázis-módosítás, fájlírás, redirect-cél vagy más protected effect előtt történjen meg.}

\diagnosticblockhu{Verifier diagnostic}{SEC-A01-005}{Egy mutation query úgy kapott objektumkulcsot, hogy nincs authorization proof az adott model/key párra.}{Töltsd be az objektumot, authorize-old, majd az authorize-olt kulcsot (vagy változatlan lokális aliasát) add át a mutationnek.}

A CRUD módosítás egy ciklust követ: olvasás, typed form input, validáció, autorizáció, tranzakciós módosítás, majd redirect és újraolvasás. A kulcskülönbség: a \emph{valid input még nem authorized input}.

\section{A request-életciklus alkalmazása CRUD módosításra}

Egy tipikus szerkesztési kérésnél a hétlépcsős modell konkrétan így olvasható:

\begin{lstlisting}
POST /notes/42
  -> parse: Form<NoteEdit>
  -> validate: title/body bounds
  -> authenticate: user session required
  -> authorize: load note 42, prove ownership/permission
  -> transaction: optimistic-lock update + audit
  -> response: typed redirect to GET
\end{lstlisting}

Három szabály a legfontosabb. A validáció soha nem bizonyít authorizationt. Az authentication soha nem bizonyít objektumtulajdont. A tranzakció előtt sikeres read soha nem helyettesíti a mutationben végrehajtott authoritative concurrency checket. Ha ezeket külön tartod, sok CRUD sérülékenység nem tud egyszerű formhibának álcázódni.

\section{Typed SQL és CRUD}

\subsection*{Lekérdezés}

\begin{lstlisting}[language=Velran]
#[query]
fn articleBySlug(db: Db, slug: Slug)
    -> Result<Article, DbError> sql {
    SELECT id, slug, title, body
    FROM articles
    WHERE slug = :slug
}
\end{lstlisting}

Az SQL szöveg statikus és compiler-owned. Felhasználói érték csak typed bindként, például \code{:slug}, kerülhet bele.

\subsection*{Mutáció transactionben}

\begin{lstlisting}[language=Velran]
#[query]
fn articleById(db: Db, id: i64)
    -> Result<Article, DbError> sql {
    SELECT id, slug, title, body
    FROM articles
    WHERE id = :id
}

#[query]
fn updateArticle(
    tx: Transaction,
    id: i64,
    title: String,
    body: String
) -> Result<(), DbError> mutates Article by id sql {
    UPDATE articles
    SET title = :title, body = :body
    WHERE id = :id
}
\end{lstlisting}

Az action betölti az objectet, létrehozza az authorization proofot, és csak ezután használja az autorizált object kulcsát a mutációhoz:

\begin{lstlisting}[language=Velran]
#[action]
fn update(
    ctx: ActionContext,
    db: Db,
    id: i64,
    title: String,
    body: String
) -> Result<Redirect, PageError> {
    let article = articleById(db, id)?;
    authorize article authenticated;
    transaction db {
        updateArticle(tx, article.id, title, body)?;
    }
    return Ok(redirect(adminArticles()));
}
\end{lstlisting}

Hétköznapi lokális tranzakciónál siker esetén commit, ismert adatbázishibánál rollback történik. A \code{Transaction} capability nem escape-elhet a blokkból. Retry-sensitive critical műveletnél viszont explicit transaction outcome-ot kell kezelni ahelyett, hogy minden hibáról rollbacket feltételeznénk; a \code{CommitUnknown} szemantikát a \ref{ch:db-deep-dive-hu}. fejezet részletezi.

\subsection*{Portable adatbázis-séma}

A V1 SQLite, PostgreSQL és MariaDB irányban is hordozható mintát céloz. \code{Date}, \code{DateTime}, \code{Uuid} és \code{Decimal} portable reprezentációja canonical text lehet.

\subsection*{Migration}

A szerver nem migrál automatikusan startupkor. Ajánlott folyamat:

\begin{lstlisting}
migrate verify
migrate apply
migrate verify
\end{lstlisting}

A migration credential legyen elkülönítve a normál application runtime credentialtől.


\section{Formok, validáció és PRG}

\subsection*{Named form schema}

\begin{lstlisting}[language=Velran]
form ArticleForm {
    title<String>
    body<String>
    published<bool>
    validate title length 2 200 body length 1 200000
}
\end{lstlisting}

Route:

\begin{lstlisting}[language=Velran]
route articleCreate POST "/admin/articles"
    form ArticleForm
    auth role Editor
    => create;
\end{lstlisting}

A handler paramétersorrendje és típusa kövesse a form mezőit.

\subsection*{Hibás beküldés}

Named formnál a hiányzó kötelező mező, a mezőszintű typed decode hiba vagy a deklarált validation hibája 422 \emph{Unprocessable Content}. A runtime a mezőértékeket escape-elve tölti vissza, mezőhibát ad és CSRF tokent illeszt a formba. Ismeretlen vagy duplikált mező viszont 400, mert az request-schema megsértése.

Ez szándékosan más, mint az inline form- vagy JSON-schema: ott a typed decode/validation hibája request-contract sértésként \code{400}. Named formot akkor használj, amikor a hibát ugyanazon felhasználói űrlapon, mezőszintű visszajelzéssel akarod kezelni.

\subsection*{Sikeres POST és PRG}

Sikeres módosítás után redirectet adj vissza; így a böngésző GET-tel tölti be a céloldalt. Ez a Post/Redirect/Get minta megszünteti a véletlen újraküldés tipikus problémáját.

\subsection*{Egyszerű teljes form-action minta}

\begin{lstlisting}[language=Velran]
#[action]
fn create(
    ctx: ActionContext,
    db: Db,
    title: String,
    body: String,
    published: bool
) -> Result<Redirect, PageError> {
    let articleSlug = slug(title);
    transaction db {
        insertArticle(tx, articleSlug, title, body)?;
    }
    return Ok(redirect(adminArticles()));
}
\end{lstlisting}

A checkbox-szerű hiányzó \code{bool} mező named form esetén \code{false}.

\subsection*{Gyakorlati döntési sorrend}

Állapotváltoztató form tervezésénél és review-jánál ezt a sorrendet használd:

\begin{enumerate}
  \item decode és validation a route/form boundaryn;
  \item meglévő rekord módosításánál töltsd be a cél objectet;
  \item hozd létre az object/permission/tenant authorization proofot;
  \item a megfelelő transaction vagy critical-operation contract alatt módosíts;
  \item typed GET route helperre redirectelj;
  \item csak azt a cache-t invalidáld, amelyet a módosítás ténylegesen elavulttá tett.
\end{enumerate}

Ez a sorrend kizár egy gyakori hibát: a jól formázott identifier vagy form önmagában nem bizonyítja, hogy a hívó módosíthatja a hivatkozott objectet.


\section{Teljes szerkesztési ciklus: validáció, authorization, konkurencia}

Létrehozó formnál gyakran elég a validáció és a tranzakció. Meglévő üzleti object szerkesztésénél azonban megjelenik egy harmadik kérdés: a felhasználónak megjelenített példány a POST beérkezésére már elavulhat. A Velran ezeket a kérdéseket külön kezeli; egy sikeres validáció nem válik automatikusan minden további lépés bizonyítékává.

Gyakorlati edit-flow:

\begin{enumerate}
  \item a request dekódolása a deklarált form-típusokra;
  \item mező- és cross-field validáció;
  \item az aktuális object betöltése és authorization proof létrehozása;
  \item a beküldött verzió ellenőrzése optimistic-locking mutációval;
  \item a módosítás és szükség esetén a business audit commitja ugyanabban a tranzakcióban;
  \item siker után redirect egy GET route-ra.
\end{enumerate}

A verzióellenőrzés magába a mutációba tartozik. Egy külön ``létezik-e még ez a verzió?'' SELECT újra race-t hozna létre az ellenőrzés és az írás között.

\begin{lstlisting}[language=Velran]
model Article {
    id: i64
    title: String
    version: i64
}

#[query]
fn updateArticle(
    tx: Transaction,
    id: i64,
    title: String,
    version: i64
) -> Result<Changed, DbError> mutates Article by id sql {
    UPDATE articles
    SET title = :title, version = version + 1
    WHERE id = :id AND version = :version
}
\end{lstlisting}

A \code{Changed} azt jelenti, hogy pontosan egy sornak kell módosulnia. Nulla érintett sor stale szerkesztést jelent és conflict lesz belőle, nem pedig valaki újabb módosításának csendes felülírása. Egynél több sor a mutáció saját invariantjának megsértése, ezért adatbázishiba.

\subsection*{Négy boundary, amit nem szabad összemosni}

\begin{longtable}{L{0.20\textwidth}L{0.31\textwidth}L{0.39\textwidth}}
\toprule
Boundary & Kérdés & Tipikus eredmény \\
\midrule
Request schema & Szerkezetileg megengedett-e a request? & ismeretlen/duplikált mező: 400 \\
Named-form validation & A felhasználó javíthatja-e ezt az értéket? & field/type/deklarált validation hiba: 422 \\
Authorization & Cselekedhet-e ez a principal ezen az objecten? & fail closed a mutáció előtt \\
Concurrency invariant & Még aktuális-e a beküldött állapot? & stale optimistic-lock verzió: 409 \\
\bottomrule
\end{longtable}

Review közben ez a szétválasztás különösen hasznos: a 422 nem authorization döntés, és a sikeres authorization sem bizonyítja, hogy egy régi böngészőfül még a legfrissebb adatbázis-állapotot tartalmazza.

\section{Adatbázis-invariantok és felhasználóbarát validáció}

Az application validation javíthatja a visszajelzést, de az authoritative uniqueness helye az adatbázis. Ha slugnak, emailnek, külső azonosítónak vagy hasonló értéknek egyedinek kell lennie, ezt adatbázis \code{UNIQUE} constraint kényszerítse ki. Preflight SELECT használható barátságosabb hibaüzenethez, de nem helyettesítheti a constraintet, mert két párhuzamos request egyszerre is átmehet a preflight ellenőrzésen.

Általánosabban is ez a minta: presentation-orientált ellenőrzés a request boundaryn, domain refinement typed/domain értékekben, authorization a védett effect előtt, race-sensitive integritás pedig az adatbázisban vagy más authoritative platform-primitívben.

\section{Mutációs review-checklist}

CRUD mutáció elfogadása előtt ellenőrizd:

\begin{itemize}
  \item a route explicit access policyval rendelkezik;
  \item a request mezők typed és bounded értékek;
  \item object-level authorization előtt betöltöd a célobjectet;
  \item SQL-be érték csak typed bindon keresztül kerül;
  \item az állapotmódosítás a megfelelő transaction/critical-operation contract alatt fut;
  \item ahol lost update számít, a konkurens szerkesztés optimistic lockingot használ;
  \item a uniqueness és más race-sensitive invariant authoritative módon van kikényszerítve;
  \item rekonstruálandó üzleti módosításnál a business audit tranzakciós;
  \item siker után PRG történik, nem mutable POST response renderelése;
  \item a cache invalidáció szűk, és csak a ténylegesen módosított adathoz kötődik.
\end{itemize}

\section{Ellenőrzött companion példák}
A fejezet copy/paste szempontból elsődleges, compilerrel ellenőrzött repository-programjai:
\begin{itemize}
  \item \path{examples/crud/app.vrn} -- CRUD és objektumszintű autorizáció;
  \item \path{examples/forms/app.vrn} -- típusos formok és validáció;
  \item \path{examples/optimistic-locking/main.vrn} -- elavult szerkesztés elleni védelem;
  \item \path{examples/prg-flash/app.vrn} -- PRG és flash állapot.
\end{itemize}

\section{Gyakorlat: tervezz szerkesztési ütközést}
\paragraph{Gyakorlási mód: támogatott.} Használd az előszóban meghatározott támogatott módszert.

Szimulálj két böngésző-sessiont, amelyek ugyanazt a rekordot nyitják meg. Az A session ment először, majd a B session elavult adatot küld vissza. Döntsd el, hol utazik a verzió, hol derül ki az eltérés, milyen HTTP eredmény születik, és mit lát ezután a felhasználó.

\section{Tipikus hiba: csak a tranzakció előtt validálunk}
Az alkalmazásszintű validáció jobb hibaüzeneteket ad, de a race-sensitive invariánsoknak továbbra is kell authoritative enforcement pont. Uniqueness, tulajdonosfüggő módosítás és optimistic concurrency akkor is maradjon helyes, ha két kérés szinte egyszerre fut.

\section{Folyamatos mintaprojekt: Secure Notes}
Használd az \path{examples/secure-notes/checkpoints/05-crud/main.vrn} compilerrel ellenőrzött snapshotot. Adj szerkesztést és törlést a Secure Notes-hoz. Legyen explicit a tulajdonosi ellenőrzés, szerkesztésnél használj optimistic lockingot, a módosítás maradjon a megfelelő tranzakciós határon belül, siker után pedig PRG-vel térj vissza.
